\chapter{روش‌شناسی ارزیابی}
\label{ch:method}

این فصل می‌گوید نتایج فصل~\ref{ch:results} \emph{چگونه} به دست آمده‌اند. تفکیک آن از فصل
نتایج عمدی است: بیشتر ادعاهای یک ارزیابی تجربی نه با خودِ اعداد، بلکه با روشی که اعداد را
تولید کرده است زیر سؤال می‌رود.

ساختار فصل چنین است: نخست بستر آزمایش و کالیبراسیون آن، سپس بارهای کاری و بازوهای مقایسه،
سپس تعریف دقیق سنجه‌ها، و سرانجام چارچوب آزمایش با بررسی‌های اعتبار آن. بخش پایانی به
دام‌هایی می‌پردازد که خودِ چارچوب در جریان ساخته‌شدن از آن‌ها آموخت \mdash{} بخشی که در بیشتر
گزارش‌ها نوشته نمی‌شود اما بیشترین ارزش عملی را دارد.

\section{سه سطح ارزیابی}
\label{sec:method-levels}

ارزیابی در سه سطح انجام شد، همان‌گونه که در پیشنهاد پروژه آمده بود:

\begin{enumerate}
  \item \textbf{آزمون واحد} برای منطق موتور سیاست، مستقل از خوشه.
  \item \textbf{آزمون یکپارچه‌سازی} برای رفتار سامانه در خوشه‌ی واقعی.
  \item \textbf{آزمون کارایی و مقایسه} در برابر سازوکار پیش‌فرض \lr{Kubernetes}.
\end{enumerate}

بیشتر این فصل و فصل بعد به سطح سوم می‌پردازد، چون ادعاهای اصلی پروژه از آنجا برمی‌آید. اما
سطح نخست از نظر روش‌شناختی جایگاه ویژه‌ای دارد: چون واسط‌ها تفکیک شده‌اند، یک حلقه‌ی کنترل
\emph{بسته} کاملاً درون یک آزمون واحد اجرا می‌شود. آنجا هیچ نویزی از شبکه، زمان‌بندی یا
تأخیر جمع‌آوری وجود ندارد، پس هر نوسانی که دیده شود قطعاً از خود سیاست است. این همان چیزی
است که ادعای بخش~\ref{sec:design-perreplica} را از یک همبستگی به یک سازوکار ارتقا می‌دهد.

\section{بستر آزمایش}
\label{sec:method-testbed}

همه‌ی اندازه‌گیری‌ها روی یک خوشه‌ی تک‌گرهی \lr{Kubernetes} انجام شد. جدول~\ref{tab:method-env}
محیط مرجع را ثبت می‌کند.

\begin{table}[H]
\centering
\small
\caption{محیط مرجع اندازه‌گیری‌ها}
\label{tab:method-env}
\begin{tabular}{|l|l|}
\hline
\textbf{مؤلفه} & \textbf{نسخه} \\
\hline
\lr{Go} & \lr{1.26.4} \\
\lr{kubectl} & \lr{v1.36.1} \\
\lr{Kubernetes} (\lr{Docker Desktop}) & \lr{v1.36.1}، تک‌گره \\
\lr{Kubernetes} (\lr{kind}) & \lr{v1.36.1}، سه‌گره \\
\lr{Helm} & \lr{v4.2.0} \\
\lr{kind} & \lr{v0.32.0} \\
\lr{k6} & \lr{v2.1.0} \\
\lr{kube-prometheus-stack} & \lr{prometheus-operator v0.92.1} \\
\lr{Prometheus Adapter} & \lr{v0.12.0} \\
تصویر پایه & \lr{golang:1.26-alpine} (ساخت)، \lr{distroless/static-debian12} (اجرا) \\
\textbf{تصویر برنامه‌ی نمونه} & \textbf{\lr{sha256:d089f870ba7a...}}، تثبیت‌شده با برچسب \\
 & \lr{sample-app:phase6-measured} \\
\hline
میزبان & \lr{macOS} (\lr{Darwin 25.5.0})، \lr{arm64} \\
ماشین مجازی & ۱۴ هسته، ۷/۷۴۹ گیگابایت حافظه، ۱۱۰ جایگاه \lr{Pod} \\
\hline
\end{tabular}
\end{table}

\lr{Prometheus} با فاصله‌ی جمع‌آوری سراسری ۱۵ ثانیه پیکربندی شد، اما برای دو
\lr{ServiceMonitor} که در اندازه‌گیری نقش دارند این فاصله ۵ ثانیه است. پنجره‌ی همه‌ی
پرس‌وجوهای \lr{rate} برابر یک دقیقه است.

سطر آخر جدول اهمیت ویژه‌ای دارد و در بخش~\ref{sec:results-image} به آن بازمی‌گردیم: تصویر
برنامه‌ی نمونه با اثرانگشت تثبیت شده است، چون یک برچسب تصویر قابل تغییر است و «همان برچسب»
تضمین «همان تصویر» نیست.

\subsection{چرا خوشه‌ی تک‌گره}
\label{sec:method-singlenode}

پیشنهاد پروژه \lr{k3s} یا \lr{kind} را نام برده بود، در حالی که توسعه‌ی روزانه روی
\lr{Kubernetes} داخلی \lr{Docker Desktop} انجام شد. این یک انحراف است و به‌جای آنکه صرفاً
ثبت شود، جبران شده است: یک پیکربندی سه‌گره‌ی \lr{kind} در مخزن قرار دارد و کل مسیر روی آن نیز
راستی‌آزمایی شده است \mdash{} سه گره آماده، هر دو تصویر بارگذاری‌شده، پخش نمونه‌ها میان دو گره کارگر،
دقیقاً یک هدف جمع‌آوری به ازای هر نمونه، و بازگرداندن یک \lr{kubectl scale} دستی توسط
مقیاس‌گذار.

با این حال، اندازه‌گیری‌های گزارش‌شده روی خوشه‌ی تک‌گره انجام شده‌اند. دلیل آن، حذف یک متغیر
مزاحم است: روی چند گره، رقابت زمان‌بندی و تفاوت بار گره‌ها به نتایج نویز اضافه می‌کند بی‌آنکه
به پرسش پژوهش ربطی داشته باشد. پیامد این انتخاب برای اعتبار نتایج در
بخش~\ref{sec:results-otherlimits} صریحاً گزارش شده است.

\section{کالیبراسیون برنامه‌ی نمونه}
\label{sec:method-calibration}

انصاف مقایسه به یک تساوی ساده وابسته است. برنامه‌ی نمونه برای هر درخواست حدود ۲ میلی‌ثانیه
پردازنده می‌سوزاند و درخواست منابع هر نمونه ۲۰۰ میلی‌کور است. پس:

\begin{equation}
\label{eq:calibration}
100~\mathrm{req/s} \times 2~\mathrm{ms} = 200~\mathrm{ms/s} = 200\mathrm{m}
\end{equation}

یعنی نمونه‌ای که در نرخ هدف خود کار می‌کند، دقیقاً ۱۰۰ درصد درخواست پردازنده‌اش را مصرف
می‌کند.

اهمیت این تساوی در آن است که دو بازو را به \emph{یک نقطه‌کار} می‌رساند. بازوی \arm{hpa-cpu}
روی \lr{averageUtilization: 100} تنظیم شده و بازوی \arm{ours-threshold} روی نرخ هدف ۱۰۰
درخواست بر ثانیه. بدون این تساوی، این دو به دو نقطه‌ی متفاوت نشانه می‌رفتند و هیچ تفاوتی
میانشان قابل انتساب نبود. تنظیم متعارف \lr{averageUtilization: 70} این تساوی را بی‌سروصدا
می‌شکست.

تساوی را اندازه گرفتیم و به استدلال اکتفا نکردیم. نخست روی میزبان و به‌صورت متوالی با
\lr{GOMAXPROCS=1}:

\begin{table}[H]
\centering
\caption{کالیبراسیون محلی برنامه‌ی نمونه}
\label{tab:method-calib-local}
\begin{tabular}{|l|r|r|}
\hline
\textbf{کمیت} & \textbf{اندازه‌گیری‌شده} & \textbf{انتظار} \\
\hline
تأخیر هر درخواست & \lr{2.096 ms} & \lr{2 ms} + سربار \lr{HTTP} \\
گذردهی تک‌هسته & \lr{477 req/s} & $\approx$\lr{500} \\
پردازنده‌ی ضمنی در \lr{100 req/s} & \lr{210m} & \lr{200m} \\
بازبینی با \lr{?cpu\_ms=8} & \lr{8.111 ms} & ۴ برابر حالت ۲ میلی‌ثانیه \\
\hline
\end{tabular}
\end{table}

سپس درون خوشه، با نگه داشتن بار روی یک نمونه و خواندن مصرف پردازنده از \lr{Prometheus}
(جدول~\ref{tab:method-calib-cluster}):

\begin{table}[H]
\centering
\caption{تأیید کالیبراسیون درون خوشه}
\label{tab:method-calib-cluster}
\begin{tabular}{|r|r|r|r|}
\hline
\textbf{بار عرضه‌شده (\lr{req/s})} & \textbf{تحویل‌شده} & \textbf{پردازنده (هسته)} & \textbf{میلی‌ثانیه بر درخواست} \\
\hline
\lr{100} & \lr{100.0} & \lr{0.232} & \lr{2.32} \\
\lr{200} & \lr{194.1} & \lr{0.431} & \lr{2.22} \\
\lr{400} & \lr{402.9} & \lr{0.862} & \lr{2.14} \\
\lr{600} & \lr{600.1} & \lr{1.271} & \lr{2.12} \\
\hline
\end{tabular}
\end{table}

مصرف اندازه‌گیری‌شده میان ۲/۱ تا ۲/۳ میلی‌ثانیه بر درخواست است و با ۲/۰۹۶ میلی‌ثانیه‌ی محلی
هم‌خوان است. در نرخ ۱۰۰ درخواست بر ثانیه مصرف ۲۳۲ میلی‌کور بود، در برابر درخواست ۲۰۰
میلی‌کوری. تساوی~\ref{eq:calibration} در خوشه نیز برقرار است.

افزون بر این، برنامه باید واقعاً \emph{محدود به پردازنده} باشد. سرویسی که به‌جای محاسبه فقط
بخوابد، زیر بار فزاینده مصرف پردازنده‌ی ثابتی نشان می‌دهد؛ آنگاه \arm{hpa-cpu} هرگز واکنش
نشان نمی‌دهد و خط مبنا به یک مترسک تبدیل می‌شود. سطر آخر
جدول~\ref{tab:method-calib-local} همین را تأیید می‌کند: چهار برابر کردن کار هر درخواست،
تأخیر را دقیقاً چهار برابر می‌کند.

\subsection{فروپاشی مشاهده‌پذیری زیر بار اضافی}
\label{sec:method-observability-collapse}

این بخش یک یافته‌ی تجربی را ثبت می‌کند که پیکربندی بستر آزمایش را تغییر داد و از نظر مفهومی
به یافته‌ی مرکزی پروژه نزدیک است.

در نخستین اجرای واقعی، بار ۱۳۰۰ درخواست بر ثانیه در برابر یک نمونه با محدودیت ۴۰۰ میلی‌کور
عرضه شد. ناوگان \emph{اصلاً مقیاس نگرفت} و در تمام مدت جهش روی یک نمونه ماند. زنجیره‌ی علت،
که از رویدادهای \lr{kubelet}، وضعیت هدف‌های \lr{Prometheus} و گزارش‌های خود مقیاس‌گذار
بازسازی شد، چنین است:

\begin{enumerate}
  \item محدودیت ۴۰۰ میلی‌کور یعنی سهمیه‌ی ۴۰ میلی‌ثانیه در هر دوره‌ی ۱۰۰ میلی‌ثانیه‌ای
  زمان‌بند، پس کانتینر زیر بار اضافی پایدار ۶۰ درصد زمان واقعی \emph{منجمد} است.
  \item نقطه‌ی \lr{/healthz} نمی‌تواند درون مهلت کاوشگر آمادگی پاسخ دهد، پس
  \lr{readyReplicas} به صفر می‌رود.
  \item همان نقطه نمی‌تواند درون مهلت کاوشگر زنده‌بودن پاسخ دهد، پس \lr{kubelet} کانتینر را
  سه بار بازراه‌اندازی می‌کند \mdash{} که بار اضافی را بدتر می‌کند و تاریخچه‌ی سنجه را نابود می‌کند.
  \item نقطه‌ی \lr{/metrics} نمی‌تواند درون مهلت جمع‌آوری پاسخ دهد، پس هدف \lr{down} می‌شود و
  \lr{rate(http\_requests\_\-total[1m])} \emph{هیچ‌چیز} برنمی‌گرداند.
  \item جمع‌آورنده نتیجه‌ی تهی را \emph{به‌درستی} صفر می‌خواند \mdash{} چون بردار تهی به‌راستی
  می‌تواند یعنی هیچ نمونه‌ای بالا نیست \mdash{} و سیاست ناوگان را روی کمینه نگه می‌دارد در حالی که
  بار بالا می‌رود.
\end{enumerate}

\begin{quote}
\textbf{یک مقیاس‌گذار که با سنجه‌های برنامه کار می‌کند، دقیقاً در لحظه‌ای که بیش از همه به آن
نیاز است نابینا می‌شود، مگر آنکه برنامه زیر بار اضافی مشاهده‌پذیر بماند.}
\end{quote}

این نقص مقیاس‌گذار نیست. رفتار آن در برابر بردار تهی عمدی و درست است. این یک ویژگی
\emph{معماری} است، و از این جهت با یافته‌ی مرکزی پروژه هم‌خانواده است: هر دو مسیرهای پس‌خوری
هستند که در آن‌ها هر مؤلفه دقیقاً طبق مشخصات خود رفتار می‌کند و مجموع، شکست می‌خورد.

سه اصلاح انجام شد. \textbf{محدودیت پردازنده از ۴۰۰ به ۱۰۰۰ میلی‌کور} \mdash{} که مهم‌ترین بود و در
انصاف مقایسه هزینه‌ای ندارد، چون ریاضیات بهره‌وری \lr{HPA} روی \emph{درخواست} پردازنده کار
می‌کند و نه روی محدودیت آن، پس تساوی~\ref{eq:calibration} دست‌نخورده می‌ماند.
\textbf{هم‌زمانی کران‌دار} با ریختن فوری درخواست‌های مازاد، تا بار عرضه‌شده در هر حالت
اندازه‌گیری‌پذیر بماند. و \textbf{کاوشگر زنده‌بودن مداراگر} با مهلت و آستانه‌ی بزرگ‌تر، چون
کاوشگر زنده‌بودن باید فرایند \emph{قفل‌شده} را تشخیص دهد و نه فرایند \emph{کُند}.

پس از این اصلاحات، همان اجرا بار کامل ۱۳۰۰ درخواست بر ثانیه را بدون هیچ بازراه‌اندازی تحویل
داد.

\section{بارهای کاری}
\label{sec:method-workloads}

چهار الگوی بار به کار رفت که هر کدام $D = 1800$ ثانیه طول می‌کشد. هر الگو یک تابع صریح از
زمان است و نه فهرستی از پله‌های دست‌نویس، چون همین توابع عیناً در شبیه‌ساز نیز به کار
می‌روند و مقایسه‌ی فصل~\ref{ch:results} به یکسان بودن منحنی‌ها وابسته است.

\begin{align}
\label{eq:pattern-spike}
\lambda_{\mathrm{spike}}(t) &=
\begin{cases}
1300 & 600 \le t < 1200 \\
300 & \text{در غیر این صورت}
\end{cases} \\[2mm]
\label{eq:pattern-diurnal}
\lambda_{\mathrm{diurnal}}(t) &= 300 + 900 \sin^2\!\left(\frac{\pi t}{D}\right) \\[2mm]
\label{eq:pattern-bursty}
\lambda_{\mathrm{bursty}}(t) &=
\begin{cases}
1200 & t \in [s, s+120) \quad s \in \{300, 800, 1300\} \\
200 & t \in [s+120, s+180) \\
400 & \text{در غیر این صورت}
\end{cases} \\[2mm]
\label{eq:pattern-ramp}
\lambda_{\mathrm{ramp}}(t) &=
\begin{cases}
200 & t < 300 \\
200 + \dfrac{1000\,(t-300)}{1200} & 300 \le t < 1500 \\
1200 & t \ge 1500
\end{cases}
\end{align}

هر الگو یک پرسش متفاوت می‌پرسد:

\textbf{\arm{spike}} یک جهش ناگهانی و بیش از چهاربرابری است که ده دقیقه دوام می‌آورد.
انگیزه‌ی پیشنهاد پروژه بر پایه‌ی همین الگو نوشته شده بود. اما \mdash{} و این یکی از نتایج منفی
گزارش‌شده در این نوشتار است \mdash{} پیش‌بینی در این الگو کاری از پیش نمی‌برد، چون هیچ
پیش‌بینی‌کننده‌ای یک پله‌ی آنی را از قبل نمی‌بیند.

\textbf{\arm{diurnal}} افت‌وخیزی نرم و سینوسی است، شبیه منحنی ترافیک روزانه. این بهترین حالت
برای پیش‌بینی است. نکته‌ی مهم آن که در فصل بعد به کار می‌آید: این تنها الگوی مجموعه است که
\emph{بی‌نهایت بار مشتق‌پذیر} است.

\textbf{\arm{bursty}} سه رگبار پیاپی است که پس از هر کدام گودالی کوتاه می‌آید. این الگو پرهیز
از نوسان را می‌آزماید. گودال، دعوتی است به مقیاس‌گیری به پایین، درست پیش از آنکه رگبار بعدی
برسد.

\textbf{\arm{ramp}} یک روند صعودی پایدار است و برای این پروژه افزوده شد. سه الگوی اصلی، بار
کاری‌ای را که پیش‌بینی واقعاً برای آن وجود دارد در بر نمی‌گرفتند: \arm{spike} و \arm{bursty}
توابع پله‌ای‌اند و \arm{diurnal} گرچه روند دارد، برمی‌گردد. \arm{ramp} یک روند یکنوای پیوسته
است، پس پیش‌بینی‌کننده چیزی واقعی برای برون‌یابی دارد. معیار روشنی به دست می‌دهد:
\emph{اگر پیش‌بینی روی \arm{ramp} برنده نشود، هیچ‌جا برنده نمی‌شود}.

\subsection{مدل بار باز، نه بسته}
\label{sec:method-openmodel}

مولد بار~\cite{k6} از سناریوی \lr{ramping-arrival-rate} استفاده می‌کند، یعنی نرخ
\emph{عرضه‌شده} را تثبیت می‌کند و نه تعداد کاربران مجازی را. این انتخاب اهمیت روش‌شناختی دارد.

در یک مدل بسته، وقتی سرویس کُند می‌شود، مولد بار خودبه‌خود کمتر می‌فرستد. یعنی بار عرضه‌شده
دقیقاً در همان لحظه‌ای افت می‌کند که مقیاس‌گذار در حال آزموده‌شدن روی واکنش به بار بالاست.
چنین چیزی همه‌ی بازوها را به یک اندازه خوش‌جلوه می‌کند و دقیقاً همان کم‌تأمین منابعی را پنهان
می‌کند که ارزیابی برای سنجش آن وجود دارد.

\subsection{گرم‌کردن}
\label{sec:method-warmup}

هر اجرا با یک مرحله‌ی گرم‌کردن آغاز می‌شود که نرخ آن، نرخ ابتدایی همان الگوست. دو دلیل دارد.
نخست آنکه \code{http\_requests\_total} تا پیش از نخستین درخواست وجود ندارد و
\lr{rate(...[1m])} پیش از گذشت یک دقیقه‌ی کامل معنایی ندارد؛ بدون گرم‌کردن، دقیقه‌ی نخست هر
اجرا صرف پر شدن خط لوله‌ی سنجه می‌شود. دوم آنکه به این ترتیب هر بازو از \emph{نقطه‌ی تعادل
خودش} وارد پنجره‌ی اندازه‌گیری می‌شود و یک شیب سرد یکسان به ابتدای همه‌ی نمودارها اضافه
نمی‌شود.

\section{بازوهای مقایسه}
\label{sec:method-arms}

جدول~\ref{tab:method-arms} نُه بازوی مقایسه را فهرست می‌کند.

\begin{table}[H]
\centering
\small
\caption{بازوهای مقایسه}
\label{tab:method-arms}
\begin{tabular}{|l|p{7.6cm}|}
\hline
\textbf{بازو} & \textbf{توضیح} \\
\hline
\arm{static} & ناوگان ثابت ۸ نمونه، بدون کنترل‌گر. تأمین‌شده برای \emph{میانگین} بار \\
\arm{static-peak} & ناوگان ثابت به اندازه‌ی $\lceil \max_t \lambda(t)/T \rceil$. تأمین‌شده برای \emph{اوج} \\
\arm{hpa-cpu} & \lr{HPA} پیش‌فرض روی مصرف پردازنده، با \lr{averageUtilization: 100} \\
\arm{hpa-custom} & \lr{HPA} روی همان سنجه‌ی سفارشی نرخ درخواست، از راه \lr{Prometheus Adapter} \\
\arm{ours-threshold} & مقیاس‌گذار این پروژه با سیاست آستانه‌ای \\
\arm{ours-predictive} & مقیاس‌گذار این پروژه با پیش‌بینی روی \textbf{بار کل} \\
\arm{ours-predictive-\-per-replica} & همان سیاست، اما پیش‌بینی روی \textbf{بار هر نمونه} \\
\arm{ours-predictive-\-nostab} & بازوی حذفی: سیاست درست، با میراگر خاموش \\
\arm{ours-predictive-\-per-replica-\-nostab} & بازوی حذفی: سیاست معیوب، با میراگر خاموش \\
\hline
\end{tabular}
\end{table}

سه نکته درباره‌ی این مجموعه اهمیت دارد.

\textbf{نخست، چرا دو خط مبنای ایستا؟} در نخستین اجراهای معتبر، بازوی واکنشی هم‌زمان بر خط
مبنای ایستا در \emph{هر دو} محور برتری داشت: ۵/۲ درصد نقض در برابر ۳۵/۵ درصد، آن هم با
نمونه-ثانیه‌ی \emph{کمتر}. چنین چیزی یک منحنی مبادله‌ی باورپذیر نیست، و علتش این بود که خط
مبنا بد تعریف شده بود. عدد ۸ تقریباً برابر نیاز \emph{میانگین} است، پس آن بازو در هر اوج
کم‌تأمین است. یک خط مبنا که در هر دو محور می‌بازد، مترسک است و داور حق دارد این را بگوید.

راه حل، نگه داشتن \emph{هر دو} سر منحنی است. \arm{static} عمداً باقی ماند، چون ظرفیت‌سنجی بر
پایه‌ی میانگین یک راهبرد واقعی است و نشان دادن شکست آن زیر بار متغیر، خودش یک نتیجه است.
\arm{static-peak} از روی تعریف الگو در زمان اجرا حساب می‌شود و نه به‌صورت عدد ثابت، تا هرگز
از بار کاری فاصله نگیرد. با این کار، ادعای جالب‌تر قابل بیان می‌شود: پرسش درست «آیا
مقیاس‌گذاری خودکار بهتر از ایستاست» نیست \mdash{} که با انتخاب یک ایستای بد قابل دستکاری است \mdash{} بلکه
«هر سیاست کجای منحنی هزینه/کیفیت می‌نشیند» است.

\textbf{دوم، چرا \arm{hpa-custom} بازوی تعیین‌کننده است.} مقایسه‌ی سیاست ما با \arm{hpa-cpu}
دو چیز را هم‌زمان تغییر می‌دهد: سیگنال (نرخ درخواست به‌جای پردازنده) و سیاست (پیش‌بینانه
به‌جای واکنشی). هر تفاوتی در نتیجه، میان این دو تغییر قابل انتساب نیست. بازوی \arm{hpa-custom}
دقیقاً \emph{همان سنجه} و \emph{همان نقطه‌ی هدف} را دارد و نزدیک‌ترین آزمون حذفی موجود در
این پروژه است.

دو تفاوت دیگر اما باقی می‌ماند و باید ثبت شود. نخست، \arm{ours-predictive} ناحیه‌ی مرده
ندارد: این سیاست شمار نمونه را مستقیماً از بار کل پیش‌بینی‌شده می‌گیرد و تابعی را صدا
می‌زند که پارامتر تلورانس نمی‌پذیرد، حال آنکه سیاست آستانه‌ای و هر دو بازوی \lr{HPA}
ناحیه‌ی مرده‌ی ۱۰ درصدی را اعمال می‌کنند؛ کنترل‌گر بی‌ناحیه‌ی مرده زودتر حرکت می‌کند، پس
بخشی از برتری زمان واکنش ممکن است از این تفاوت بیاید و نه از پیش‌بینی
(بخش~\ref{sec:conclusion-future}). دوم، \arm{hpa-custom} سیگنال را از راه
\lr{prometheus-adapter} می‌گیرد و دو پرش اضافی دارد؛ اینکه این مسیر گران است از خود داده
پیداست، چون \arm{hpa-cpu} و \arm{hpa-custom} یک سیاست یکسان دارند اما روی \arm{diurnal}
زمان واکنششان ۵۵/۰ در برابر ۱۱۲/۵ ثانیه است. فاصله‌ی باقی‌مانده همچنان بزرگ است، اما ادعا
باید با این دو قید بیان شود.

\textbf{سوم، بازوهای \arm{-nostab}.} این دو بازو با \lr{stabilization\_window\_seconds: 0}
اجرا می‌شوند و همه چیز دیگرشان یکسان است. وجود آن‌ها پاسخ مستقیم به پرسش پژوهشی دوم است و در
بخش~\ref{sec:results-stabilizer} نتیجه‌ی اصلی پروژه از آن‌ها می‌آید.

هر بازو در هر لحظه \emph{تنها} فعال است. دو کنترل‌گر که هم‌زمان روی یک \lr{Deployment}
بنویسند با هم می‌جنگند و رد به‌جامانده تفسیرناپذیر است؛ چارچوب آزمایش این را به‌صورت یک گیت
اجباری تضمین می‌کند.

\section{تعریف سنجه‌ها}
\label{sec:method-metrics}

پیشنهاد پروژه چهار سنجه را نام برده بود: زمان واکنش به جهش بار، تأخیر پاسخ به کاربر، میزان
کم‌تأمین و بیش‌تأمین منابع، و مصرف کلی منابع. سنجه‌ی پنجمی نیز افزوده شد که در پیشنهاد نبود و
در عمل تعیین‌کننده از آب درآمد: \textbf{تعداد تغییر جهت}.

تعریف‌ها بر پایه‌ی «نیاز مرجع» بیان می‌شوند، که از خودِ تعریف الگو می‌آید و نه از بار
تحویل‌شده:

\begin{equation}
\label{eq:needed}
n(t) = \left\lceil \frac{\lambda(t)}{T} \right\rceil
\end{equation}

استفاده از تعریف الگو به‌جای بار اندازه‌گیری‌شده عمدی است: اگر مخرج مقایسه خودش تحت تأثیر
رفتار بازو باشد، بازویی که ناوگانش فرو می‌پاشد به‌طور خودکار با معیار آسان‌تری سنجیده می‌شود.

این قاعده بر $U$ و $O$ حاکم است، که از $n(t)$ می‌آیند. اما $V$ از سری
\emph{اندازه‌گیری‌شده}‌ی بار هر نمونه محاسبه می‌شود و نه از الگو. پیامدش در جهت
محافظه‌کاری است: وقتی ناوگانی فرو می‌پاشد، بار سرویس‌شده کاهش می‌یابد و $V$ آن بازو
\emph{کم‌برآورد} می‌شود؛ پس ارقام بالای نقض کف مقدارند و نه سقف آن
(بخش~\ref{sec:results-slameaning}).

\begin{description}
  \item[نقض توافق سطح خدمت ($V$)] درصد نمونه‌های زمانی که در آن‌ها بار هر نمونه از خط
  $1.1\,T = 110$ درخواست بر ثانیه فراتر رفته است. بخش~\ref{sec:results-slameaning} روشن
  می‌کند این سنجه دقیقاً چه چیزی را می‌سنجد و چه چیزی را نمی‌سنجد.

  \item[هزینه ($R$)] نمونه-ثانیه، یعنی $\int_{t_0}^{t_1} r(t)\,dt$ روی پنجره‌ی اندازه‌گیری.

  \item[کم‌تأمین ($U$)] $\int \max(0,\, n(t) - r(t))\,dt$.

  \item[بیش‌تأمین ($O$)] $\int \max(0,\, r(t) - n(t))\,dt$.

  \item[زمان واکنش] برای هر پله‌ی صعودی در $n(t)$، فاصله‌ی زمانی تا لحظه‌ای که
  $r(t) \ge n(t)$ شود. میانه‌ی این مقادیر گزارش می‌شود. اگر ناوگان تا پایان پنجره هرگز به
  عدد لازم نرسد، این کمیت \emph{تعریف‌نشده} است \mdash{} و این خودش یک نتیجه است و نه یک عدد گمشده.

  \item[تغییر جهت ($N_{\mathrm{rev}}$)] شمار دفعاتی که علامت تغییر \lr{spec.replicas} نسبت به
  تغییر پیشین عوض می‌شود. مستقیم‌ترین سنجه‌ی نوسان همین است.

  \item[مصرف پردازنده] هسته-ثانیه، از سنجه‌های \lr{cAdvisor}.

  \item[تأخیر] صدک ۹۵ از بافه‌ی \code{http\_request\_duration\_seconds} خود برنامه.
\end{description}

دو نکته‌ی پیاده‌سازی درباره‌ی این تعریف‌ها ارزش ثبت دارد.

\textbf{کنش‌های مقیاس‌گذاری روی \lr{spec.replicas} شمرده می‌شوند و نه روی
\lr{status.readyReplicas}.} \lr{spec} همان کمیتی است که کنترل‌گر واقعاً می‌نویسد؛ شمردن روی
\lr{ready} پایان یافتن راه‌اندازی نمونه‌ها را هم می‌شمرد، که واکنش \emph{خوشه} است و نه تصمیم
\emph{سیاست}.

\textbf{صدک ۹۹ تأخیر از مولد بار گرفته نمی‌شود.} خروجی خلاصه‌ی \lr{k6} تنها صدک‌های ۹۰ و ۹۵ را
حمل می‌کند؛ آستانه‌ی صدک ۹۹ را ارزیابی می‌کند اما مقدارش را دور می‌ریزد. بنابراین همه‌ی
ارقام تأخیر از بافه‌ی خود برنامه می‌آید و دید سمت‌کاربر \lr{k6} تنها به‌عنوان بررسی متقابل
نگه داشته می‌شود.

\section{چارچوب آزمایش}
\label{sec:method-harness}

چارچوب آزمایش پیش از اجرای هر آزمایشی ساخته شد، و این ترتیب عمدی بود. راه دیگر \mdash{} اجرای بیست
آزمایش دستی و پردازش بعدی آن‌ها \mdash{} همان راهی است که در آن آدم در آزمایش بیستم می‌فهمد آزمایش
سوم نامعتبر بوده، بی‌آنکه بتواند بگوید کدام‌یک از بقیه همان نقص را دارند.

دو برنامه کار را انجام می‌دهند و مرز میانشان سخت است: یکی \emph{یک اجرا} را انجام می‌دهد و
داده‌ی خام را ذخیره می‌کند، دیگری همه‌ی سنجه‌ها را از روی داده‌ی خام \emph{مشتق} می‌کند.
هیچ سنجه‌ای در حین اجرا محاسبه نمی‌شود و هیچ چیزی در حین تحلیل اندازه‌گیری نمی‌شود.

اهمیت این تفکیک عملی است: چون سری‌های خام نگه داشته می‌شوند، تعریفی از یک سنجه که بعداً
نادرست از آب دربیاید قابل اصلاح است و همه‌ی اجراهای گذشته بدون بازگشت به خوشه دوباره
امتیازدهی می‌شوند. در جریان این پروژه دقیقاً همین اتفاق دو بار افتاد.

\subsection{یک اجرا چیست}
\label{sec:method-run}

یک اجرا برابر با «اجرای مولد بار» نیست. مراحل آن چنین است:

\begin{enumerate}
  \item \textbf{برچیدن.} همه‌ی کنترل‌گرهای احتمالی برداشته می‌شوند.
  \item \textbf{اعمال دقیقاً یک بازو.}
  \item \textbf{بازنشانی ناوگان} به وضعیت آغازین.
  \item \textbf{گرم‌کردن} با نرخ ابتدایی الگو.
  \item \textbf{اندازه‌گیری} به مدت ۱۸۰۰ ثانیه.
  \item \textbf{نشست} برای ثبت رفتار پس از پایان بار.
  \item \textbf{ضبط ۱۷ سری زمانی} از \lr{Prometheus}.
  \item \textbf{بررسی اعتبار} و نوشتن فراداده‌ی اجرا.
\end{enumerate}

هر مرحله‌ای جز «اندازه‌گیری»، یک یافته‌ی تجربی است که به یک گیت تبدیل شده است. دو مورد از
همه پرهزینه‌تر بودند:

\textbf{بازخوانی \lr{ConfigMap} زنده پس از تعویض سیاست.} ویرایش یک \lr{ConfigMap} نمونه را
بازراه‌اندازی نمی‌کند، و یک \lr{ConfigMap} سوارشده تنها در دوره‌ی همگام‌سازی \lr{kubelet}
به‌روز می‌شود. بدون یک \lr{rollout restart} صریح به‌علاوه‌ی بازخوانی، سیاست \emph{بازوی قبلی}
زیر نام بازوی جدید به کار خود ادامه می‌دهد. آنگاه هر عددی به سیاست اشتباه نسبت داده می‌شود و
هیچ‌جا هیچ‌چیز مشکوکی دیده نمی‌شود.

\textbf{گرم‌کردن}، به دلایلی که در بخش~\ref{sec:method-warmup} گفته شد.

\subsection{اعتبار و دود: دو محور جداگانه}
\label{sec:method-validity}

دو صفت جداگانه در فراداده‌ی هر اجرا ثبت می‌شود و یکی‌گرفتن آن‌ها نخستین خطای طراحی چارچوب
بود:

\begin{itemize}
  \item \textbf{نامعتبر} یعنی مکانیک اجرا شکسته و هیچ‌چیز قابل بازیابی نیست.
  \item \textbf{دود}\LTRfootnote{Smoke} یعنی مقیاس زمانی برابر یک نیست. در چنین اجرایی
  ممکن است مکانیک بی‌نقص باشد و اعداد همچنان هیچ جایی در فصل نتایج نداشته باشند.
\end{itemize}

بررسی‌های اعتباری که به‌طور خودکار انجام می‌شوند:

\begin{itemize}
  \item نام سیاست فعال از \lr{ConfigMap} زنده بازخوانی و مقابله می‌شود.
  \item طول پنجره‌ی اندازه‌گیری باید در بازه‌ی مورد انتظار بماند. در عمل ۱۸۰۰ تا ۱۸۰۲ ثانیه بود.
  \item تعداد هدف‌های جمع‌آوری باید با تعداد نمونه‌های آماده برابر باشد.
  \item نرخ \code{http\_requests\_total} پس از گرم‌کردن باید ناصفر باشد.
  \item شمارش \code{http\_reqs} خود مولد بار باید دست‌کم ۹۵ درصد انتگرال الگو باشد.
  \item تعداد بازراه‌اندازی‌های برنامه‌ی نمونه باید صفر باشد.
  \item در بازوهای ایستا، \lr{spec.replicas} باید ثابت \emph{و برابر مقدار اعلام‌شده} بماند.
  \item اثرانگشت تصویر برنامه‌ی نمونه باید با تصویر تثبیت‌شده یکی باشد.
\end{itemize}

\textbf{یک شکاف شناخته‌شده در همین فهرست.}
\label{sec:validity-gate-gap}
دروازه‌ی تحویل بار آنچه \lr{k6} \emph{فرستاده} را می‌سنجد و نه آنچه به ناوگان \emph{رسیده}.
این دو تنها وقتی از هم جدا می‌شوند که نقص مسیریابی یا فروپاشی ناوگان در کار باشد، یعنی
دقیقاً همان جایی که به دروازه نیاز است؛ نقص بخش~\ref{sec:results-loaddefect} از همین شکاف
عبور کرد. شکل درست این بررسی، مقابله‌ی \code{sum(rate(http\_requests\_total[1m]))} سمت
ناوگان با انتگرال الگوست، که آنجا دستی انجام شد اما به دروازه‌های خودکار افزوده نشد.

\subsection{چرا مقیاس زمانی باید یک باشد}
\label{sec:method-timescale}

فشرده‌سازی محور زمان وسوسه‌انگیز است: یک اجرای ۳۰ دقیقه‌ای را می‌توان در ۳ دقیقه انجام داد و
یک جاروب کامل به‌جای سه ساعت، بیست دقیقه طول می‌کشد. این کار در این پروژه ممنوع است و دلیلش
تجربی است.

یک اجرای دودی با مقیاس زمانی ۲۰ روی \arm{ramp} این را مستقیماً نشان داد: ناوگان در پایان
شیب به ۱۱ نمونه رسید و هرگز به ۱۲ لازم نرسید، و نرخ اندازه‌گیری‌شده در اوج حدود ۹۱۰ درخواست
بر ثانیه بود در برابر ۱۲۰۰ عرضه‌شده. هر دو انحراف، مصنوع محض محور فشرده‌اند: یک پنجره‌ی
\lr{rate} یک‌دقیقه‌ای در این مقیاس ۱۲۰۰ ثانیه‌ی \emph{الگو} را پوشش می‌دهد، پس سیگنال بار
تا حد ناشناختگی هموار می‌شود؛ و نمونه‌ها همان ۳۰ ثانیه برای آماده‌شدن وقت می‌برند در حالی که
بار کاری بیست برابر زودتر تمام می‌شود، پس ناوگان به‌هیچ‌وجه نمی‌تواند آن را دنبال کند.
هیچ‌کدام چیزی درباره‌ی سیاست نمی‌گویند.

نکته‌ی مهم‌تر آن است که نسبت میان «سرعت حرکت بار» و «زمان آماده‌شدن نمونه» دقیقاً همان پدیده‌ای
است که این پروژه مطالعه می‌کند. فشرده‌سازی زمان، متغیر مستقل را تغییر می‌دهد.

پیامد عملی: یک بازو حدود ۳۴ دقیقه و یک جاروب کامل حدود سه ساعت طول می‌کشد.

\subsection{دستگاه اندازه‌گیری هم باید کالیبره شود}
\label{sec:method-instrument}

سه نقص در جریان ساخت چارچوب پیدا شد که همگی در \emph{کد بررسی‌کننده} بودند و نه در سامانه‌ی
تحت آزمون. هر سه نیز در جهت \emph{رد کردن اجراهای سالم} سوگیری داشتند، که بدترین حالت ممکن
است چون تا وقتی یک روز کار خوشه را نخورده باشد دیده نمی‌شود.

\textbf{شمارش بازراه‌اندازی می‌توانست منفی شود.} گیت، مجموع شمارنده‌ها را در دو لحظه مقایسه
می‌کرد، اما آن مجموع روی \emph{مجموعه‌ی فعلی} نمونه‌هاست: هر بازویی که مقیاس به پایین بدهد
نمونه‌هایی را حذف می‌کند، شمارنده‌هایشان از مجموع خارج می‌شود، و اختلاف منفی می‌شود. یعنی گیت
روی هر بازوی خودکار اجراهای سالم را رد می‌کرد، \emph{و هرچه سیاست بهتر مقیاس به پایین می‌داد
بیشتر}.

\textbf{گیتی که برای تشخیص کنترل‌گر بیگانه ساخته شده بود، خودِ چارچوب را تشخیص داد.} پنجره‌ی
ضبط به پیش از آغاز اندازه‌گیری بازمی‌گشت و بنابراین لحظه‌ی برپاسازی بازو را هم در بر می‌گرفت.
اکنون این پرسش روی پنجره‌ی \emph{اندازه‌گیری} پرسیده می‌شود، که همان پنجره‌ای است که ادعای
اعتبار درباره‌ی آن است.

\textbf{گیتی که به دلیل اشتباه شلیک می‌کند از نبودِ گیت فقط اندکی بهتر است.} یک اجرا که بر
اثر خواب میزبان نابود شده بود، با پیام «مولد بار تکرارها را انداخته است» رد شد. پیام، مؤلفه‌ی
اشتباهی را متهم می‌کرد و نفر بعدی را به سراغ اشکال‌زدایی چیز دیگری می‌فرستاد. «زمان
ازدست‌رفته» اکنون دلیل نامعتبری مستقل خودش را دارد.

جمع‌بندی این سه مورد یک اصل است: \textbf{چارچوب آزمایش خودش یک دستگاه اندازه‌گیری است، و
دستگاهی که هرگز روی ورودی معلوم خوانده نشده، کالیبره نشده است.} به همین دلیل، هر بازوی جدید
پیش از صرف ۳۴ دقیقه، یک بار به‌صورت دودی آزموده می‌شود.

\section{تکرارپذیری و قدرت تفکیک}
\label{sec:method-repeatability}

هیچ مقایسه‌ای معنا ندارد مگر بدانیم پراکندگی یک بازو در برابر \emph{خودش} چقدر است. برای همین
چند بازو سه بار به‌صورت مستقل اجرا شد. مقادیر به‌جای میانگین‌گیری تک‌تک گزارش می‌شوند، چون در
$n = 3$ میانگین بیش از آنکه آشکار کند پنهان می‌کند.

قاعده‌ای که از این کار به دست می‌آید در سراسر فصل~\ref{ch:results} رعایت شده است:

\begin{quote}
\textbf{تفاوتی میان دو بازو که از پراکندگی یک بازو در برابر خودش کوچک‌تر باشد، تفاوت نیست.}
\end{quote}

اعداد این پراکندگی در بخش~\ref{sec:results-repeatability} آمده است. نکته‌ی روش‌شناختی مهم آن
است که تکرارها بر پایه‌ی \emph{شرایط اجرا} گروه‌بندی می‌شوند و نه فقط بر پایه‌ی نام بازو:
اجراهای دو سوی یک تغییر در بستر، دو آزمایش‌اند و نه دو تکرار، و درهم‌آمیختنشان فاصله‌ی میان
آن دو را به‌عنوان «نویز این بازو» گزارش می‌کند. چون همین پراکندگی، قدرت تفکیک کل ارزیابی را
تعیین می‌کند، یک پراکندگی متورم بی‌سروصدا یافته‌های واقعی را رد صلاحیت می‌کند.

\section{جمع‌بندی فصل}
\label{sec:method-summary}

روش‌شناسی این ارزیابی سه ویژگی دارد که ادعاهای فصل بعد به آن‌ها وابسته است.

نخست، \textbf{انصاف مقایسه اندازه‌گیری شده است و فرض نشده}: کالیبراسیون
برابری~\ref{eq:calibration} هم روی میزبان و هم درون خوشه تأیید شده، و بازوی \arm{hpa-custom}
همان سیگنال و همان نقطه‌ی هدف بازوی ما را دارد.

دوم، \textbf{اعتبار هر اجرا به‌صورت خودکار و پیش از ورود به تحلیل بررسی می‌شود}، و
«نامعتبر» از «دودی» جدا شمرده می‌شود.

سوم، \textbf{قدرت تفکیک اندازه‌گیری پیش از هر مقایسه‌ای تعیین شده است}، و هر ادعایی که از آن
باریک‌تر باشد در فصل بعد به‌عنوان نشانه گزارش می‌شود و نه نتیجه.

با این حال \mdash{} و این نکته‌ای است که بخش~\ref{sec:results-validity} به آن اختصاص دارد \mdash{} در
جریان کار نقصی کشف شد که از تمام این بررسی‌ها گذشت و همچنان معیوب بود. گزارش صادقانه‌ی آن
بخشی از نتیجه است.
